\histitem{Fredholm index and perturbations}

The notion of index appears in 1920 in a work by Fritz Noether on integral equations. In his lecture at the 1904 International Congress, Hilbert considered a problem posed by Riemann in potential theory: given a bounded open domain \(\Omega\) of the plane whose boundary is a smooth simple closed curve \(\Gamma\) , as well as three functions \(a\) , \(b\) , \(f\) on \(\Gamma\) , the problem is to find a function \(z: \overline{\Omega} \to \mathbf{C}\) , holomorphic on \(\Omega\) and continuous on \(\overline{\Omega}\) , satisfying

\[
a \mathcal {R} (z) + b \mathcal {I} (z) = f \quad \text { on } \Gamma .
\]

Parametrizing \(\Gamma\) by a real number \(s \in [0, 2\pi]\) , Hilbert showed that the problem reduces to solving an integral equation “with singular kernel” for \(\varphi = \mathcal{R}(z)\) , namely

\[
a (s) \varphi (s) + \int_ {0} ^ {2 \pi} \mathrm{K} (s, t) \varphi (t) d t = f (s).\tag{4}
\]

Here the kernel \(\mathrm{K}(s, t)\) takes the form \(b(s) \cot\left(\frac{t-s}{2}\right) + \mathrm{A}(s, t)\) for a continuous function A. It is therefore singular along the diagonal, and the above integral is to be taken in the sense of the Cauchy principal value.

Noether {[}53{]} observes that, unlike integral equations satisfying the Fredholm alternative, in the case of a kernel K as above and a nonzero right-hand side \(f\), it is possible that (4) admits families of nontrivial solutions. He shows that the space of solutions is governed by the integer

\[
n = \frac {1}{2 \pi} \int_ {\Gamma} d (\log (a - i b)).\tag{5}
\]

If \(n < 0\), there is no nontrivial solution, and if \(n \geqslant 0\), equation (4) admits a family of solutions depending on 2n parameters. Noether uses the term "index" (Index) for the integer n and recognizes in it a winding number, a notion classical since the end of the nineteenth century in the study of functions of a complex variable.

Later works would go on to study examples of equations requiring an extension of Fredholm theory. But it would take more than twenty years before the notion of a Fredholm map was systematically extracted, at the same time as the study of perturbations by compact operators matured.

In 1909, H. Weyl {[}81{]} had shown that if two quadratic forms bounded on \(\ell_{\mathbf{C}}^{2}(\mathbf{N})\) have a completely continuous difference, then their essential spectra coincide (cf. III, p. 89, th. 2). Moreover, Riesz had noted in his 1916 memoir, naturally without the modern language, that the compact operators on a Banach space E form a two-sided ideal of \(\mathcal{L}(\mathrm{E})\). It seems, however, that the path indicated by these two remarks was not explored until the 1940s. In 1941, J. Calkin, motivated by the work of J. von Neumann on operator algebras (which we shall have occasion to discuss shortly, see p. 537), pointed out the interest of studying congruences modulo this ideal in the framework of Hilbert spaces {[}10{]}. He shows how the structure of the algebra that now bears his name (cf. III, p. 75) makes it possible to recover Weyl's result simply.

Meanwhile, around 1942, J. Dieudonné {[}15{]} became interested in the idea of studying perturbations of a continuous linear map \(u\) between normed spaces, in the case where the perturbation is “small” in the sense of the operator norm. He restricts himself to the case where \(u\) is a strict morphism whose kernel is finite-dimensional or of finite codimension; he proves that the same then holds for any “small perturbation” of \(u\) , and shows that in the case of Fredholm maps ( \(cf. n^{\circ} 2\) from III, p. 40), the index is locally constant (th. 1 of III, p. 58), as are several of the results presented in § 4 of III, p. 55.

These two ideas converge around 1950, when F. Atkinson {[}3{]}, I. Gohberg {[}26{]}, S. Mikhlin {[}46{]} , and B. Yood {[}93{]} study perturbations by compact operators. They explicitly make the connection with the Noether index and bring out the notions of Fredholm map {[}3, p. 8{]} and of Riesz map {[}93, §5{]}. The latter are the subject of systematic research in the following decade, where the results presented in Chapter III take on more or less their current form.

The usual framework for these theories is that of Banach spaces; however, various applications motivate their generalization to broader classes of topological vector spaces. Thus, the case of Fréchet spaces appears naturally in 1954 in works of H. Cartan and J-P. Serre {[}13{]} which demonstrate the finiteness of the cohomology of a compact complex analytic manifold with coefficients in a coherent sheaf, the proof invoking deformation results proved by L. Schwartz {[}66{]} (cf. th. 2 of III, p. 73).

The invariance of the index under deformation is perfectly illuminated, in the case of formula (5), if one observes that the right-hand side (winding number) is invariant under homotopy. Such topological expressions of the index would soon become famous for Fredholm applications arising from differential geometry. If D is a differential operator on a compact \(C^{\infty}\) manifold X (cf. VAR, §14), and if D is “elliptic” (cf.{[}2, §1{]}), then an extension of D to suitable Sobolev spaces defines a Fredholm map. The search for a formula giving a topological expression for its index using the “characteristic classes” of X, suggested among others by the work of I. Gelfand around 1960 {[}23, vol. I, p. 65{]}, will lead in 1963 to the “index theorem” of M. Atiyah and I. Singer {[}2{]}. This latter would give rise to extraordinary developments at the confluence of analysis, topology, and differential geometry. We cannot describe here the vast fields they opened, which, still very fertile today, have not ceased to draw nourishment from Fredholm theory.

